\section{Ablation Studies}
\label{sec:ablation-studies}

To better understand the scalability and robustness of harness evolution, we conduct ablation studies examining two critical factors. 
First, we investigate how expanding the available resource budget impacts final performance. Second, we analyze whether the evolved harnesses successfully generalize across different underlying policy models. 

\subsection{Effect of Evolution Budget}

To evaluate the scaling behavior of harness evolution, we track the performance of Qwen3.7-Max and GLM-5.2 across three different budget constraints: 24 hours, 10 iterations with 500 steps; 36 hours, 15 iterations with 750 steps; and 48 hours, 20 iterations with 1000 steps. As illustrated in Figure~\ref{fig:budget-ablation}, both models demonstrate a clear positive scaling trend. Expanding the budget from 24 to 48 hours yields consistent, monotonic improvements in both the Overall Score and the Anytime Validation score.

Notably, GLM-5.2 exhibits a sharp performance trajectory, achieving a steep climb in its Anytime Validation metric during the initial expansion to 36 hours before plateauing. Conversely, Qwen3.7-Max demonstrates a steadier, linear growth across all measured constraints. This positive scaling trend indicates that despite the previously noted limitations in budget utilization, allocating greater computational resources for iterative exploration can still provide meaningful performance dividends, ultimately enabling the synthesis of more refined agent harnesses.

\subsection{Effect of the Policy Model}

Our main experiments fix DeepSeek-V4-Flash as the policy model. To verify whether the harness evolution capability is tied to a specific policy model, we swap it to Qwen3.6-35B-A3B and GLM-5.2, and re-evaluate the evolution process using Qwen3.7-Max and GLM-5.2 as evolvers. 

% 需要润色一下
Table~\ref{tab:policy-model-ablation} demonstrates that harness evolution is highly robust to policy model changes. Regardless of whether the policy agent is based on Qwen, DeepSeek, or GLM, the evolved harnesses consistently achieve massive improvements over their respective CodeAct baselines. For instance, when utilizing Qwen3.6-35B-A3B as the policy model, the initial Overall score starts at a modest 13.9. Yet, both evolvers successfully optimize the harness to reach scores of 27.9 and 29.2. A similar trend emerges with the highly capable GLM-5.2 policy model, where the GLM-5.2 evolver lifts the baseline score from 38.0 to an impressive 48.4. These consistent cross-policy gains demonstrate that the evolvers are genuinely synthesizing generalizable reasoning and tool-use structures rather than merely overfitting to the idiosyncratic flaws of a single target model.




\definecolor{PolicyBlue}{HTML}{E3F5FB}

\begin{figure}[htbp]
\centering
\begin{minipage}[t]{0.47\linewidth}
    \vspace{0pt}
    \centering
    \includegraphics[width=\linewidth]{Figures/budget_scaling.pdf}
    \caption{Effect of evolution budget.}
    \label{fig:budget-ablation}
\end{minipage}
\hfill
\begin{minipage}[t]{0.51\linewidth}
    \vspace{0pt}
    \centering
    \scriptsize
    \setlength{\tabcolsep}{1.8pt}
    \renewcommand{\arraystretch}{1.03}
    \begin{tabular}{@{}lrrrrr@{}}
    \toprule
    \textbf{Evolver}
    & \textbf{Search}
    & \textbf{Office}
    & \textbf{General}
    & \textbf{Overall}
    & \textbf{ATV} \\
    \midrule

    \rowcolor{PolicyBlue}
    \multicolumn{6}{c}{\textit{Policy: Qwen3.6-35B-A3B}} \\
    Baseline
    & 2.7 & 14.2 & 35.9 & 13.9 & -- \\
    Qwen3.7-Max
    & 12.5 & 33.0 & \textbf{48.4} & 27.9 & 29.2 \\
    GLM-5.2
    & \textbf{16.4} & \textbf{34.0} & 45.3
    & \textbf{29.2} & \textbf{30.8} \\
    \midrule

    \rowcolor{PolicyBlue}
    \multicolumn{6}{c}{\textit{Policy: DeepSeek-V4-Flash}} \\
    Baseline
    & 11.7 & 38.4 & 48.4 & 29.7 & -- \\
    Qwen3.7-Max
    & 36.3 & 37.8 & \textbf{59.4} & 41.5 & 49.3 \\
    GLM-5.2
    & \textbf{45.4} & \textbf{39.2} & 48.4
    & \textbf{43.5} & \textbf{51.0} \\
    \midrule

    \rowcolor{PolicyBlue}
    \multicolumn{6}{c}{\textit{Policy: GLM-5.2}} \\
    Baseline
    & 18.0 & 40.2 & 73.4 & 38.0 & -- \\
    Qwen3.7-Max
    & \textbf{38.3} & 45.1 & 46.9 & 42.7 & 46.7 \\
    GLM-5.2
    & 35.6 & \textbf{45.2} & \textbf{80.3}
    & \textbf{48.4} & \textbf{50.4} \\

    \bottomrule
    \end{tabular}
    \captionof{table}{Effect of the policy model on harness evolution. ATV
    denotes AnytimeVal; baseline is the CodeAct loop.}
    \label{tab:policy-model-ablation}
\end{minipage}
\end{figure}

% TODO(results): Fill Table~\ref{tab:policy-model-ablation} and analyze ranking
% stability and cross-policy transfer of the evolved harness structures.


% \subsection{Effect of Task Diversity}

% To isolate the value of multi-domain evidence, we compare a general harness evolved on all five sources with harnesses evolved only on search, office, or general tasks. GLM-5.2 is held fixed as the evolver. This ablation tests whether diverse train tasks produce broadly transferable components or whether domain-specific evolution is more effective within its own task family.

% \begin{table}[htbp]
% \centering
% \small
% \caption{Effect of task diversity with GLM-5.2 as the evolver.}
% \label{tab:task-diversity-ablation}
% \resizebox{\columnwidth}{!}{
% \begin{tabular}{l l r}
% \toprule
% Harness & Evolution train data & Eval Gain $\uparrow$ \\
% \midrule
% Seed                  & --                   & 0.0000 \\
% Mixed general harness & All five sources     & -- \\
% Search harness         & BrowseComp + HLE     & -- \\
% Office harness         & GDPval + APEX-Agents & -- \\
% General harness        & Claw-Eval            & -- \\
% \bottomrule
% \end{tabular}
% }
% \end{table}

% TODO(results): Define the common evaluation surface for Table~\ref{tab:task-diversity-ablation}
% explicitly once the domain runs are finalized, then report transfer gaps.
